\section{Exercícios da Lista 1}

Os exercícios desta lista foram selecionados do livro
\citetitle{zill2001equacoes}, de Dennis G. Zill e Michael R. Cullen, Terceira
Edição \cite{zill2001equacoes}. São eles: o exercício \texttt{29} da página
\texttt{59} e o exercício \texttt{23} da página \texttt{77}.

\question{
  (x^2 + xy + y^2)\, dx - xy\, dy=0.
}

\answer{
\textbf{Classificação}

\begin{itemize}
\item \textbf{Ordem:} 1ª ordem.
\item \textbf{Grau:} 1º grau.
\item \textbf{Linearidade:} não linear.
\item \textbf{Separabilidade:} não separável diretamente.
\item \textbf{Homogeneidade:} homogênea, pois

$$
  \frac{dy}{dx}
  =
  \frac{x^2+xy+y^2}{xy}
  =
  \frac{x}{y}+1+\frac{y}{x}
  =
  F\left(\frac{y}{x}\right).
$$

\end{itemize}

\textbf{Resolução}

Fazendo

$$
  y=vx,
  \qquad
  \frac{dy}{dx}=v+x\frac{dv}{dx},
$$

temos

\begin{align*}
  v+x\frac{dv}{dx}
  &=\frac{1}{v}+1+v\
  x\frac{dv}{dx}
  &=\frac{1+v}{v}.
\end{align*}

Separando as variáveis,

$$
  \frac{v}{1+v}\,dv=\frac{dx}{x}.
$$

Como

$$
  \frac{v}{1+v}
  =
  1-\frac{1}{1+v},
$$

segue que

\begin{align*}
  \int\left(1-\frac{1}{1+v}\right)dv
  &=\int\frac{dx}{x}\
  v-\ln|1+v|
  &=\ln|x|+C.
\end{align*}

Substituindo $v=\dfrac{y}{x}$,

$$
  \frac{y}{x}
  -\ln\left|1+\frac{y}{x}\right|
  =
  \ln|x|+C.
$$

Como

$$
  \ln\left|1+\frac{y}{x}\right|
  =
  \ln|x+y|-\ln|x|,
$$

obtemos

$$
  \frac{y}{x}-\ln|x+y|=C.
$$

Assim,

$$
  \ln|x+y|
  =
  \frac{y}{x}-C,
$$

e, portanto,

$$
  \boxed{x+y=Ce^{y/x}}.
$$

\textbf{Solução adicional}

Durante a separação, dividimos por $1+v$. Logo, verificamos

$$
  1+v=0
  \quad\Longrightarrow\quad
  v=-1.
$$

Como $y=vx$,

$$
  y=-x.
$$

Portanto, as soluções são

$$
 \boxed{x+y=Ce^{y/x}}
$$

e

$$
 \boxed{y=-x}.
$$

}


\question{
  \cos^2{x} \sin{x}\, dy + (y \cos^3{x} - 1)\, dx = 0.
}

\answer{
\textbf{Classificação:}

\begin{itemize}
\item \textbf{Ordem:} 1ª ordem, pois a maior derivada é $\dfrac{dy}{dx}$.
\item \textbf{Grau:} 1º grau, pois $\dfrac{dy}{dx}$ aparece na primeira
potência.
\item \textbf{Linearidade:} linear, pois pode ser escrita na forma

$$
  \frac{dy}{dx}+P(x)y=Q(x).
$$

\item \textbf{Separabilidade:} não é separável diretamente.
\item \textbf{Homogeneidade:} não homogênea, pois possui o termo independente
$-1$.
\end{itemize}

Dividindo a equação por $dx$ e por $\cos^2x\sin x$, obtemos

\begin{align*}
  \frac{dy}{dx}
  +
  \frac{y\cos^3x}{\cos^2x\sin x}
  &=
  \frac{1}{\cos^2x\sin x}\
  \frac{dy}{dx}+y\frac{\cos x}{\sin x}
  &=
  \frac{1}{\cos^2x\sin x}.
\end{align*}

Logo,

$$
  \frac{dy}{dx}+y\cot x
  =
  \sec^2x\csc x.
$$

Como a equação é linear, utilizamos o \textbf{fator integrante}

$$
  \mu(x)=e^{\int\cot x\,dx}
  =e^{\ln|\sin x|}
  =\sin x.
$$

Multiplicando toda a equação por $\sin x$,

\begin{align*}
  \sin x\frac{dy}{dx}+y\sin x\cot x
  &=\sin x\sec^2x\csc x\
  \sin x\frac{dy}{dx}+y\cos x
  &=\sec^2x.
\end{align*}

O lado esquerdo é a derivada do produto:

$$
  \frac{d}{dx}(y\sin x)=\sec^2x.
$$

Integrando,

\begin{align*}
  \int\frac{d}{dx}(y\sin x),dx
  &=\int\sec^2x,dx\
  y\sin x&=\tan x+C.
\end{align*}

Portanto,

$$
  y=\frac{\tan x+C}{\sin x}.
$$

Como

$$
  \frac{\tan x}{\sin x}
  =
  \frac{\sin x}{\cos x\sin x}
  =
  \sec x,
$$

obtemos

$$
 \boxed{y=\sec x+C\csc x}.
$$
}
